% !TEX root = Tesi_Corsi_Leonardo_656276.tex
\chapter{Introduzione}
\label{ch:introduzione}

Un array sistolico è una griglia bidimensionale di unità di calcolo identiche
che comunicano esclusivamente con i vicini immediati: i dati attraversano la
griglia in maniera sincrona, con passi scanditi dal ciclo di clock, e ogni unità, 
a ogni ciclo, elabora ciò che riceve dalle unità di calcolo adiacenti.

Il modo più semplice di realizzare un array di questo tipo è collegare le
unità con canali privi di un meccanismo di sincronizzazione esplicito: il
produttore scrive il dato nel canale quando il proprio calcolo lo produce, il
consumatore lo legge quando il proprio calcolo lo richiede, e nessuno dei due
segnala all'altro lo stato del canale. La correttezza di questo schema
dipende allora per intero dalla schedulazione temporale con cui l'host
alimenta l'array, che deve conoscere in anticipo a quale ciclo ciascuna unità
si aspetta ciascun dato. Se lo scarto fra l'istante di produzione e quello di
consumo non coincide con quello previsto dalla schedulazione, una cella
riceve un operando associato a un'iterazione diversa da quella attesa: il
calcolo procede regolarmente, ma su dati non corrispondenti, e produce un
risultato numericamente definito ma errato, senza che alcun componente rilevi
l'anomalia.
\clearpage
La tesi adotta un principio elementare per rimuovere questa dipendenza dalla
schedulazione: impostare canali sincroni che usano \emph{ready bit} in grado di
modificare l'eventuale flusso di dati nel canale. Ciascun lato del canale possiede
un registro da un bit e commuta soltanto il proprio, il produttore quando il dato è
pronto e il consumatore quando lo ha prelevato; è il confronto fra i due registri,
e non il valore dell'uno o dell'altro, a dire se il canale è pieno o vuoto. Nessuno
dei due lati deve quindi conoscere in anticipo il comportamento temporale
dell'altro, e il risultato del calcolo smette di dipendere dalla velocità relativa
dei nodi.

La verifica di questo principio è oggetto della costruzione di
\textbf{NESO}\footnote{Il nome corrisponde alle quattro direzioni cardinali
(nord, est, sud, ovest), le sole lungo le quali una cella può comunicare.}, un
simulatore di array bidimensionale di processori RISC-V, corredato dalla
modellazione in SystemVerilog del meccanismo di canale sincrono e dai programmi che ne
verificano il comportamento.

\section{Il lavoro in sintesi}
\label{sec:sintesi}

NESO simula una griglia $R \times C$ di celle. Ogni cella è un processore
RISC-V completo, con i suoi trentadue registri, il suo \emph{program counter} e
la sua memoria privata; le celle non condividono memoria e comunicano solo con i
quattro vicini ortogonali attraverso canali monodirezionali di profondità uno.
Il tempo è globale e discreto: a ogni ciclo di clock ogni cella viva esegue
esattamente una istruzione.

Il punto di partenza è l'interprete a singolo processore RISC-V sviluppato in una
tesi precedente~\cite{TesiStella}. A partire da quello sono stati realizzati:

\begin{itemize}
    \item il \textbf{canale} e il suo protocollo, in cui tutto lo stato di
          sincronizzazione è costituito da due contatori a un bit e dai
          comparatori che li confrontano;
    \item quattro \textbf{istruzioni di comunicazione} (\texttt{IN},
          \texttt{OUT}, \texttt{ISRDY}, \texttt{SETRDY}) aggiunte all'insieme
          RV32I usando lo spazio di codifica \emph{custom-0} che la
          specifica RISC-V lascia libero, senza modificare l'assemblatore;
    \item la \textbf{griglia}: topologia, cablaggio dei vicini, clock globale con
          aggiornamento a due fasi e ingresso/uscita sul perimetro;
    \item nove \textbf{programmi di prova} in assembly, con i programmi di
          verifica in C che li eseguono e ne controllano il risultato con delle
          asserzioni: dal protocollo minimo fra due celle alla moltiplicazione
          di matrici e allo stencil di Jacobi;
    \item il \textbf{modello RTL} del canale e dell'interfaccia di cella in
          SystemVerilog, verificato con due \emph{testbench} e sintetizzato per
          contare quante porte costa il protocollo;
    \item la \textbf{parallelizzazione} del simulatore con OpenMP, insieme alla relativa analisi
          di scalabilità
\end{itemize}
Il simulatore, i nove programmi in assembly e i programmi di verifica in C che
li misurano sono stati sviluppati interamente nel corso di questa tesi: i dati
riportati nel capitolo~\ref{ch:risultati} sono il loro risultato diretto.

\section{Organizzazione del lavoro}
\label{sec:organizzazione}

Il capitolo~\ref{ch:strumenti} presenta gli strumenti impiegati: l'architettura
RISC-V e il sottoinsieme RV32I, i linguaggi usati per le tre viste del sistema
(C per il simulatore, assembly per i programmi, SystemVerilog per il modello
hardware) e OpenMP. Il capitolo~\ref{ch:progetto} descrive il progetto
logico, ovvero l'architettura dell'array e la struttura fisica di una cella e
di un canale, e discute i due piani su cui il lavoro è parallelo: il calcolo
distribuito sulle celle e l'esecuzione del simulatore stesso. Il
capitolo~\ref{ch:implementazione} tratta i dettagli più singificativi dell'
implementazione: la codifica delle quattro istruzioni, il ciclo di clock a due 
fasi su cui si fonda la sincronizzazione, il protocollo del canale con l'invariante 
che ne garantisce la correttezza, lo stato dei registri di ogni cella 
e la corrispondenza con il modello in SystemVerilog. Il capitolo~\ref{ch:risultati} raccoglie i
risultati: cosa dimostra ciascun programma di prova, il confronto con il caso
privo del \emph{ready bit}, il costo in cicli in funzione della forma
della griglia e le prestazioni del simulatore.\ifconclusioni{} Il
capitolo~\ref{ch:conclusioni} riassume il lavoro svolto e l'esperienza di
sviluppo, e indica gli sviluppi futuri.\fi
